\section{The exact Type I seed interface}
\label{sec:cubic-typei-seed}

\begin{lemma}[AC-K29-seed-shape: the cubic tuple is not a scalar adjoint tuple]
\label{AC-K29-seed-shape}
Let $H$ be a group with an involutive automorphism $c$.
Let $\xi,\gamma:H\to k^\times$ be characters over a field,
with
\[
 \xi^3=1,\quad \xi\ne1,\quad \xi^c=\xi^{-1},
                         \qquad \gamma^c=\gamma.
                                                               \tag{SI.1}
\]
For $j=1,2$, no character $u:H\to k^\times$ and no
permutation transform the multiset
\[
                  u\{\xi^j\gamma,1,\xi^j\gamma^{-1}\}
                                                               \tag{SI.2}
\]
into $\{\theta,1,\theta^{-1}\}$ for a character $\theta$.
This applies to the cubic residual tuple at every
$\ell\ne3$, and also to its order-nine scalar normalization.
\end{lemma}
\begin{proof}
For three characters $a,b,d$, a scalar multiple of their
multiset can be $\{\theta,1,\theta^{-1}\}$ only if one
of the three characters is the geometric mean of the other
two: choosing the entry which becomes $1$ gives respectively
\[
                         bd=a^2,\qquad ad=b^2,
                                  \qquad ab=d^2.              \tag{SI.3}
\]
Take $a=\xi^j\gamma$, $b=1$, $d=\xi^j\gamma^{-1}$.
The middle alternative gives $\xi^{2j}=1$, impossible
for a nontrivial cubic character.  The first gives
$\gamma^3=\xi^{-j}$, and the third gives
$\gamma^3=\xi^j$.  Apply $c$ to either identity and use
(SI.1).  Its right side is inverted and its left side is
unchanged, again forcing $\xi^{2j}=1$.  This proves all
cases, independently of the scalar $u$.
Reduction preserves the order of $\chi$ at $\ell\ne3$,
and a rational character restricted to $G_K$ is fixed by
complex conjugation.  Thus (SI.1) holds for the residual
characters under consideration.  The order-nine scalar
normalization is among the arbitrary $u$ already excluded.
\end{proof}

\begin{remark}[AC-K29-TypeI-source: the generic source layer and the missing construction]
\label{AC-K29-TypeI-source}
The following is an audit of the actual pinned source
interfaces, rather than an assertion that cubic Type I
automorphy is false.  In this paragraph all filenames are
relative to
\path{src/Development/AutomorphicArithmetic/AutomorphicVarieties/ModularityArgument/ThresholdArguments/DenseThresholdCases/}
at commit \texttt{7cbf7a0c32fca1badf05931879a37c3f8a2b03b4}.

\begin{enumerate}
\item The literal \texttt{AutBAlt} in
\path{AutomorphyUnramifiedCoefficients.lean}, lines 16--18,
quantifies points of the old $D.\mathrm{Ring}$, where
$D=\mathrm{Dat}(\Omega,K,c,e,\psi)$ and $\psi$ is a
rational Dirichlet character.  Its Type II companion is
the pointwise input used successfully in AC-K29-typeII.
The Type I assembly in \path{AutEnd3Alt.lean},
\texttt{autEnd3\_alt}, retains the old \texttt{OrdPt}
and \texttt{GalLink} throughout its reduction to
\texttt{AutPtAlt}.  It does not itself replace the residual
tuple.

\item A genuinely more general layer is present.
\path{IsSeed.lean}, lines 25--34, defines
\texttt{IsSeedAlt} for arbitrary tuples of characters
$\psi_i,\tau_i$ on a CM field, with continuity,
irreducibility, a symmetric conjugate-dual pairing,
the residual sum, specified away-from-$\ell$ type,
ordinary local flags, and an \emph{already automorphic}
scalar twist.  Likewise
\path{PointTheoremBaseObject.lean}, lines 26--49,
quantifies arbitrary \texttt{Resid} data.
Its point clause requires \texttt{IsAutCsdA} before
concluding $J\subseteq\ker x$.
The complete assembly
\path{THidaX.lean}, \texttt{hida\_x\_of\_printed\_hida},
produces a Hecke ideal $J$ for arbitrary residual data.
The record \texttt{UnitaryHeckeData} in
\path{UnitaryHeckeData.lean} has no nonzero localized-module
field.  More decisively, the exact definition
\texttt{HSet.IsJ} in
\path{HeckeAlgebras/IsTheta.lean}, lines 46--52
(relative here to \path{AutomorphicVarieties/}), explicitly
allows the alternative
\[
          M_{K}=0,\qquad J=B,                             \tag{SI.4}
\]
before its nonzero Hecke-algebra branch.  Thus existence
of this ideal alone does not assert nontrivial automorphic
support for the cubic datum.  An actual automorphic point
to which the point clause applies excludes (SI.4).

\item The supplied construction of a rank-three seed is
\path{AdjointLiftFromEisensteinSeedOfHyps.lean},
\texttt{adjoint\_lift\_from\_eisenstein\_of\_propagation},
whose full proof starts at line 184.  It constructs a
rational two-dimensional ordinary Eisenstein-congruent
representation $\rho$, takes $\operatorname{ad}^0\rho$,
and checks the tuple $(\Psi,1,\Psi^{-1})$.
The proof explicitly uses the identity
$\operatorname{ad}^0(\rho\otimes\nu)
=\operatorname{ad}^0\rho$.
AC-K29-seed-shape shows that neither a scalar twist of
this output nor a permutation provides the cubic residual
tuple.  This is a statement about that specific seed
construction, not about all possible automorphic seeds.

\item Propagation must also be checked for the changed datum.
The exact statement \texttt{SeedPtOfAlt} in
\path{PrConc.lean}, lines 19--34, and the complete applied
proof \texttt{main}/\texttt{seedPtOf\_alt} in
\path{PrSeedPtOf.lean}, lines 152--259, still instantiate
the tuples as powers of a rational $\Psi$ through
\texttt{CReal} and \texttt{GalLink}.
The conclusion \texttt{AutPtAlt} in
\path{PrDesc.lean}, lines 16--21, has the same frame
restriction.  The generic determinant-pair argument
inside this proof accepts componentwise residual
characters, but that fact alone does not verify all
frame, auxiliary-place, and component-propagation
hypotheses for the cubic tuple.
\end{enumerate}

The remaining Type I target is consequently concrete:
construct an actual ordinary conjugate-self-dual cuspidal
seed with the cubic residual sum and its prescribed local
types, and verify the changed frame and propagation
interfaces that carry its automorphic support to every
actual Type I point used by the finite extraction.
The higher-level generic Hida statements can be retained
as source inputs once their hypotheses are verified.
This audit neither assumes that a suitable seed cannot
exist nor promotes the existence of a generic Hecke ideal
to the required pointwise automorphy theorem.
It also leaves the \texttt{Adm3U\_s} level and no-link
conditions in AC-K29-budget explicit.  No new Lean run or
independent verification of the imported automorphy
literature is claimed.
\end{remark}

\paragraph{Geometric dictionary.}
The three residual characters specify a decomposed finite
local system; scalar twisting tensors it by a rank-one
local system.  AC-K29-seed-shape is an obstruction to the
specified adjoint construction of this local system.
The formal deformation ring records lifts with their
local conditions, while the Hecke ideal records the support
provided by the automorphic module.  Its allowed zero
module branch explains why existence of the ideal does not
itself supply an automorphic point.  The finite-cover
terminology reuses the audited SGA~1, Expos\'e~V,
\S\S4 and~7 dictionary only; neither the deformation
conditions nor nonzero Hecke support is attributed to
that categorical equivalence.
